\section{功与功率}

功描述力在空间位移上的累积效应。不能把功简单理解成“力的大小乘以路程”，因为只有力在位移方向上的分量参与做功。功率则描述做功的快慢。

\subsection{恒力做功}

若物体在恒力 $\boldsymbol F$ 作用下发生位移 $\boldsymbol s$，力与位移夹角为 $\theta$，则该力做功为
\[
    W=\boldsymbol F\cdot\boldsymbol s=Fs\cos\theta .
\]
因此：

\begin{enumerate}
    \item 当 $0\leq \theta<90^\circ$ 时，$W>0$，力做正功；
    \item 当 $\theta=90^\circ$ 时，$W=0$，力不做功；
    \item 当 $90^\circ<\theta\leq 180^\circ$ 时，$W<0$，力做负功。
\end{enumerate}

功是标量。功的正负表示能量转移方向，不表示空间方向。

\subsection{变力做功}

若力随位置变化，或物体沿曲线运动，则应把路径分成许多微小位移元 $\mathrm d\boldsymbol r$。在每一小段上，元功为
\[
    \mathrm dW=\boldsymbol F\cdot\mathrm d\boldsymbol r .
\]
从起点 $a$ 到终点 $b$ 的总功为
\[
    \boxed{
    W=\int_a^b\boldsymbol F\cdot\mathrm d\boldsymbol r
    } .
\]
在直角坐标系中可写为
\[
    W=\int_a^b
    \left(
    F_x\,\mathrm dx+F_y\,\mathrm dy+F_z\,\mathrm dz
    \right).
\]
计算变力做功时，应沿实际路径积分，不能随意取始末力的平均值乘总位移。

\subsection{合力做功}

若物体受到多个力
\[
    \boldsymbol F_1,\boldsymbol F_2,\cdots,\boldsymbol F_n,
\]
则合力为
\[
    \boldsymbol F_{\mathrm{net}}=\sum_{i=1}^{n}\boldsymbol F_i .
\]
合力做功为
\[
    W_{\mathrm{net}}
    =
    \int \boldsymbol F_{\mathrm{net}}\cdot\mathrm d\boldsymbol r
    =
    \sum_{i=1}^{n}
    \int \boldsymbol F_i\cdot\mathrm d\boldsymbol r
    =
    \sum_{i=1}^{n}W_i .
\]
所以，合力做功等于各分力做功的代数和。

\subsection{功率}

平均功率定义为
\[
    \bar P=\frac{\Delta W}{\Delta t}.
\]
瞬时功率为
\[
    P=\frac{\mathrm dW}{\mathrm dt}.
\]
由 $\mathrm dW=\boldsymbol F\cdot\mathrm d\boldsymbol r$ 得
\[
    \boxed{
    P=\boldsymbol F\cdot\boldsymbol v
    } .
\]
若力与速度夹角为 $\theta$，则
\[
    P=Fv\cos\theta .
\]
功率单位为瓦特：
\[
    1\,\mathrm W=1\,\mathrm{J/s}.
\]

\subsection{功的相对性}

功与参考系有关。因为位移 $\mathrm d\boldsymbol r$ 和速度 $\boldsymbol v$ 都依赖参考系，所以同一个力在不同惯性系中做功和功率一般不同。与此同时，动能也依赖参考系；只要在同一参考系内计算，功与动能之间仍满足动能定理。

\subsection{例题：变力沿直线做功}

质点沿 $x$ 轴从 $x=0$ 运动到 $x=a$。作用在质点上的力为
\[
    F_x=kx,\qquad F_y=F_z=0,
\]
其中 $k$ 为常量。求该力做的功。

\textbf{解析：} 因为质点沿 $x$ 轴运动，
\[
    \mathrm d\boldsymbol r=\mathrm dx\,\boldsymbol e_x.
\]
元功为
\[
    \mathrm dW=\boldsymbol F\cdot\mathrm d\boldsymbol r=kx\,\mathrm dx.
\]
故
\[
    W=\int_0^a kx\,\mathrm dx=\frac{1}{2}ka^2.
\]
因此
\[
    \boxed{W=\frac{1}{2}ka^2}.
\]
本题体现了变力做功必须积分，不能简单套用恒力做功公式。

\subsection{常用公式}

\[
    W=\boldsymbol F\cdot\boldsymbol s=Fs\cos\theta
\]
\[
    \boxed{
    W=\int_a^b\boldsymbol F\cdot\mathrm d\boldsymbol r
    =
    \int_a^b
    \left(
    F_x\,\mathrm dx+F_y\,\mathrm dy+F_z\,\mathrm dz
    \right)
    }
\]
\[
    \boxed{
    P=\frac{\mathrm dW}{\mathrm dt}=\boldsymbol F\cdot\boldsymbol v
    }
\]
